% genere par scripts/rechnung_retraite.py (quellen_tex) - ne pas modifier a la main
\begin{description}
\item[\texttt{abgeltungsteuer\_satz}] \phantomsection\label{src:abgeltungsteuer_satz} §32d Abs. 1 EStG. \url{https://www.gesetze-im-internet.de/estg/__32d.html}, consulté le 29.09.2026 (source primaire).
\item[\texttt{basiszins\_2026}] \phantomsection\label{src:basiszins_2026} BMF-Schreiben vom 13.01.2026, Basiszins zur Berechnung der Vorabpauschale zum 2. Januar 2026, IV C 1 - S 1980/00230/012/001. \url{https://www.bundesfinanzministerium.de/Content/DE/Downloads/BMF_Schreiben/Steuerarten/Investmentsteuer/2026-01-13-basiszins-berechnung-vorabpauschale.html}, consulté le 29.09.2026 (source secondaire). \textit{Remarque de collecte~: Zugriff auf das BMF-Schreiben selbst lieferte keinen lesbaren Zahlenwert; Wert per WebFetch aus dem Stollfuss- und Ecovis-Bericht über das gleiche Schreiben bestaetigt (Basiszins 3,20 \%). Sekundaerbestaetigung, siehe LUECKEN.md.}
\item[\texttt{durchschnittsentgelt\_2026}] \phantomsection\label{src:durchschnittsentgelt_2026} §3 Abs. 2 Sozialversicherungsrechengroessen-Verordnung 2026 (SVBezGrV 2026). \url{https://www.gesetze-im-internet.de/svbezgrv_2026/__3.html}, consulté le 29.09.2026 (source primaire).
\item[\texttt{einstiegsgehalt\_brutto}] \phantomsection\label{src:einstiegsgehalt_brutto} StepStone Gehaltsreport 2026, Medianwert Ingenieurinnen/Ingenieure mit unter einem Jahr Berufserfahrung (Datenbasis Jan. 2022 - Nov. 2025). \url{https://www.stepstone.de/e-recruiting/hr-wissen/gehalt/gehaltsreport-deutschlands-gehaelter-im-fokus/}, consulté le 29.09.2026 (source secondaire). \textit{Remarque de collecte~: Sekundaerquelle (Personaldienstleister-Report), keine unternehmensspezifische Jungheinrich/IG-Metall-Kueste-Eingruppierung gefunden, siehe LUECKEN.md}
\item[\texttt{est\_tarif\_2026}] \phantomsection\label{src:est_tarif_2026} §32a Abs. 1 EStG, Veranlagungszeitraum 2026. \url{https://www.gesetze-im-internet.de/estg/__32a.html}, consulté le 29.09.2026 (source primaire). \textit{Remarque de collecte~: Jede Zeile: [obere Grenze der Zone (EUR, null = letzte Zone), Typ, a, b, c]. Typ "null": ESt = 0. Typ "y"/"z": mit t = (x - untere Grenze der Zone)/10000, ESt = floor((a*t + b)*t + c). Typ "lin": ESt = floor(a*x + b), b bereits mit Vorzeichen (negativ) fuer die direkte Addition, quivalent zur Gesetzesschreibweise ESt = a*x - |b|. Werte gegen zwei unabhaengige Sekundaerquellen (finanz-tools.de, rechnercheck.de) kreuzgeprueft, identisch; Vorzeichen- und Typkonvention 2026-09-30 an salaire.est\_32a angeglichen (vorher progression1/progression2/linear42/linear45 mit positivem b, c=null in den linearen Zonen).}
\item[\texttt{etf\_welt\_rendite\_div}] \phantomsection\label{src:etf_welt_rendite_div} Repli explicite (Annahme), aucune valeur exploitable trouvee sur la fiche justETF du plus grand ETF MSCI World. \url{https://www.justetf.com/en/etf-profile.html?isin=IE00B4L5Y983}, consulté le 30.09.2026 (source secondaire). \textit{Remarque de collecte~: Le plus grand ETF MSCI World par encours (iShares Core MSCI World UCITS ETF USD Acc, IE00B4L5Y983, EUR 129 840 m au 31.08.2026 selon la fiche justETF) est THESAURISANT: aucun rendement de distribution n'y est affiche (verifie par WebFetch le 30.09.2026, page consultee integralement). A titre de comparaison seulement (pas retenu comme valeur, car ce n'est pas le plus grand ETF): la part distribuante iShares MSCI World UCITS ETF (Dist), IE00B0M62Q58 (https://www.justetf.com/en/etf-profile.html?isin=IE00B0M62Q58, encours EUR 8 303 m), affiche 'Current dividend yield 0,85 \%' et 'Dividend yield (1 an) 1,01 \%', tous deux sensiblement sous le repli de 1,8 \% retenu ici. Repli 0,018 conserve tel quel (decision du controleur, tache 9), ecart documente pour revision de jour. Voir LUECKEN.md.}
\item[\texttt{gehaltssteigerung\_real}] \phantomsection\label{src:gehaltssteigerung_real} Abgeleitet aus StepStone Gehaltsreport 2026 (Querschnitt nach Berufserfahrung: Einstieg 59.250 EUR vs. 25 Jahre Berufserfahrung 93.250 EUR, impliziert ca. 1,8 \%/Jahr Wachstumsrate). \url{https://www.stepstone.de/e-recruiting/hr-wissen/gehalt/gehaltsreport-deutschlands-gehaelter-im-fokus/}, consulté le 29.09.2026 (source secondaire). \textit{Remarque de collecte~: Kein direkter Destatis-Verdiensterhebungs-Zeitreihenwert zur uber die Karriere gemittelten realen Steigerung gefunden; Querschnittsdaten nach Erfahrungsjahren als Naeherung verwendet, siehe LUECKEN.md}
\item[\texttt{inflation\_2022\_de}] \phantomsection\label{src:inflation_2022_de} Statistisches Bundesamt (Destatis), Pressemitteilung Nr. 022 vom 17.01.2023, 'Inflationsrate im Jahr 2022 bei +7,9 \%'. \url{https://www.destatis.de/DE/Presse/Pressemitteilungen/2023/01/PD23_022_611.html}, consulté le 29.09.2026 (source primaire).
\item[\texttt{inflation\_destatis\_2019\_2025}] \phantomsection\label{src:inflation_destatis_2019_2025} Statistisches Bundesamt (Destatis), Pressemitteilungen annuelles 'Inflationsrate im Jahr JJJJ': PD20\_019\_611 (2019), PD21\_025\_611 (2020), PD22\_025\_611 (2021), PD25\_020\_611 (citant retrospectivement 2022 a +6,9\% et 2023 a +5,9\%), PD24\_020\_611 (2023), PD25\_020\_611 (2024), PD26\_019\_611 KORREKTUR (2025). \url{https://www.destatis.de/DE/Presse/Pressemitteilungen/2025/01/PD25_020_611.html}, consulté le 30.09.2026 (source primaire). \textit{Remarque de collecte~: Serie pour fig\_inflation\_2022 (chapitre 7, tache 9 partie B), sourcee par avance ici. ATTENTION ecart de revision: la premiere publication provisoire de janvier 2023 (PD23\_022\_611, deja citee ici sous inflation\_2022\_de) annoncait +7,9\% pour 2022; les publications retrospectives 2024/2025 (PD24\_020\_611, PD25\_020\_611) citent uniformement +6,9\% pour 2022 en comparaison des annees suivantes, probablement lie au changement de base du VPI (nouvelle base 2020=100 introduite en 2024). Serie ci-dessus prise entierement dans les publications retrospectives coherentes entre elles (2021/2022/2023/2024 cites ensemble dans PD25\_020\_611), pas melangee avec le chiffre provisoire de 2023. inflation\_2022\_de (7,9\%) laisse inchange ailleurs dans le depot (hors perimetre de cette tache); ecart signale dans LUECKEN.md pour arbitrage de jour.}
\item[\texttt{inflation\_ziel}] \phantomsection\label{src:inflation_ziel} EZB, geldpolitische Strategie, Preisstabilitaetsziel von 2 \% (HVPI). \url{https://www.ecb.europa.eu/mopo/strategy/pricestab/html/index.de.html}, consulté le 29.09.2026 (source primaire).
\item[\texttt{kv\_bbg\_monat\_2026}] \phantomsection\label{src:kv_bbg_monat_2026} GKV-Spitzenverband-Faktenblatt 'Rechengroessen und Grenzwerte im Versicherungs- und Beitragsrecht fuer das Jahr 2026', Beitragsbemessungsgrenzen Kranken- und Pflegeversicherung, S. 1. \url{https://www.gkv-spitzenverband.de/media/dokumente/presse/zahlen_und_grafiken/20260101_Faktenblatt_Rechengroessen_Beitragsrecht.pdf}, consulté le 29.09.2026 (source primaire).
\item[\texttt{kv\_mindestbemessung\_monat\_2026}] \phantomsection\label{src:kv_mindestbemessung_monat_2026} GKV-Spitzenverband-Faktenblatt 'Rechengroessen und Grenzwerte im Versicherungs- und Beitragsrecht fuer das Jahr 2026', Freiwillige Versicherung in der Krankenversicherung, Mindestbemessungsgrundlage, S. 1. \url{https://www.gkv-spitzenverband.de/media/dokumente/presse/zahlen_und_grafiken/20260101_Faktenblatt_Rechengroessen_Beitragsrecht.pdf}, consulté le 29.09.2026 (source primaire).
\item[\texttt{kv\_satz\_ermaessigt}] \phantomsection\label{src:kv_satz_ermaessigt} §243 SGB V; bestaetigt im GKV-Spitzenverband-Faktenblatt 'Rechengroessen und Grenzwerte im Versicherungs- und Beitragsrecht fuer das Jahr 2026' (1. Januar 2026), Beitragssaetze, S. 1. \url{https://www.gkv-spitzenverband.de/media/dokumente/presse/zahlen_und_grafiken/20260101_Faktenblatt_Rechengroessen_Beitragsrecht.pdf}, consulté le 29.09.2026 (source primaire). \textit{Remarque de collecte~: PDF gespeichert unter refs/gkv-spitzenverband-rechengroessen-2026.pdf}
\item[\texttt{kv\_schwellen\_dynamisierung}] \phantomsection\label{src:kv_schwellen_dynamisierung} §223 Abs. 4 SGB V: 'Die Beitragsbemessungsgrenze im Jahr 2027 entspricht der um 3 600 Euro erhoehten Jahresarbeitsentgeltgrenze nach § 6 Absatz 7 [SGB V] ... Die Bundesregierung setzt die Beitragsbemessungsgrenze in der Rechtsverordnung nach § 160 des Sechsten Buches gesondert fest.' §6 Abs. 6 Satz 2 SGB V (Jahresarbeitsentgeltgrenze, via Abs. 7 referenziert): 'Sie aendert sich zum 1. Januar eines jeden Jahres in dem Verhaeltnis, in dem die Bruttoloehne und -gehaelter je Arbeitnehmer ... im vergangenen Kalenderjahr zu den entsprechenden Bruttoloehnen und -gehaeltern im vorvergangenen Kalenderjahr stehen', d.h. eine jaehrliche Lohn-Indexierung, keine feste nominale Grenze. Die Mindestbemessungsgrundlage fuer freiwillig Versicherte ist ein Bruchteil der Bezugsgroesse (§18 SGB IV): 'das Durchschnittsentgelt der gesetzlichen Rentenversicherung im vorvergangenen Kalenderjahr, aufgerundet auf den naechsthoeheren, durch 420 teilbaren Betrag', ebenfalls jaehrlich per Sozialversicherungs-Rechengroessenverordnung festgesetzt.. \url{https://www.gesetze-im-internet.de/sgb_5/__223.html}, consulté le 30.09.2026 (source primaire). \textit{Remarque de collecte~: Weitere Fundstellen: §6 SGB V (https://www.gesetze-im-internet.de/sgb\_5/\_\_6.html, Abs. 6-7) und §18 SGB IV (https://www.gesetze-im-internet.de/sgb\_4/\_\_18.html), alle per curl am 2026-09-30 wortgetreu abgerufen. Modellkonvention (scripts/projection.py, \_saetze\_real): da Loehne mindestens so stark wie Preise wachsen und das reale Lohnwachstum im Modell separat (salaire.py) abgebildet wird, werden Mindestbemessung und BBG in scripts/projection.py KONSTANT in Euros 2026 (real) gehalten (konservativ-neutrale Wahl, keine explizite Lohnwachstumsannahme fuer die Schwellen selbst) - anders als der Sparerpauschbetrag (§20 Abs. 9 EStG), der ein fester NOMINALER Betrag ist und sich nur per Gesetzesaenderung veraendert, daher im Modell real dekumuliert (deflatiert) wird. Korrigiert eine fruehere Fehlimplementierung, die Mindestbemessung/BBG faelschlich wie den Sparerpauschbetrag deflatierte (Maximalbeitrag faellt dadurch nominal-eingefroren von ca. 1226 auf 513 EUR/Monat bis 2071, was jede Dividendenstrategie zu guenstig aussehen liess).}
\item[\texttt{kv\_teilfreistellung\_etf}] \phantomsection\label{src:kv_teilfreistellung_etf} GKV-Spitzenverband, 'Katalog von Einnahmen und deren beitragsrechtliche Bewertung nach § 240 SGB V', Stand 26.05.2026, Zeile 'Investmenterträge': '§ 16 InvStG | ja, unter Beruecksichtigung der §§ 20 und 56 Abs. 6 InvStG'. Die separate Zeile 'Kapitalvermoegen, Einkuenfte aus -' (individuelle Aktiendividenden, kein Fondsvehikel) traegt keinen solchen Verweis auf §20 InvStG., 15 von 30, Zeile 'Investmenterträge'. \url{https://www.gkv-spitzenverband.de/media/dokumente/krankenversicherung_1/grundprinzipien_1/finanzierung/beitragsbemessung/2026-05-26_Katalog_Einnahmen_beitragsrechtliche_Bewertung_240_SGB_V_BF.pdf}, consulté le 30.09.2026 (source primaire). \textit{Remarque de collecte~: Text per pdftotext -layout selbst extrahiert und gegenkontrolliert (nicht nur ueber eine KI-Zusammenfassung des Fetch-Tools uebernommen, da eine erste WebSearch-Zusammenfassung und eine erste WebFetch-Zusammenfassung sich widersprachen). Gilt nur fuer Fondsertraege (§16 InvStG, deckt sowohl Ausschuettung als auch Vorabpauschale ab); individuelle Aktiendividenden (Poche maison) erhalten keine Teilfreistellung, weder steuerlich noch in der GKV-Basis.}
\item[\texttt{kv\_zusatzbeitrag\_2026}] \phantomsection\label{src:kv_zusatzbeitrag_2026} GKV-Spitzenverband-Faktenblatt 'Rechengroessen und Grenzwerte im Versicherungs- und Beitragsrecht fuer das Jahr 2026', Beitragssaetze, Beitragssaetze, S. 1. \url{https://www.gkv-spitzenverband.de/media/dokumente/presse/zahlen_und_grafiken/20260101_Faktenblatt_Rechengroessen_Beitragsrecht.pdf}, consulté le 29.09.2026 (source primaire). \textit{Remarque de collecte~: PDF gespeichert unter refs/gkv-spitzenverband-rechengroessen-2026.pdf; vom Bundesgesundheitsministerium fuer 2026 festgelegter Schaetzerkreis-Wert}
\item[\texttt{pv\_satz\_kinderlos\_2026}] \phantomsection\label{src:pv_satz_kinderlos_2026} GKV-Spitzenverband-Faktenblatt 'Rechengroessen und Grenzwerte im Versicherungs- und Beitragsrecht fuer das Jahr 2026', Beitragssaetze (3,6 \% + 0,6 \% Zuschlag fuer Kinderlose), Beitragssaetze, S. 1. \url{https://www.gkv-spitzenverband.de/media/dokumente/presse/zahlen_und_grafiken/20260101_Faktenblatt_Rechengroessen_Beitragsrecht.pdf}, consulté le 29.09.2026 (source primaire). \textit{Remarque de collecte~: Fuer freiwillig Versicherte traegt der Versicherte den vollen Satz allein, daher hier 0,042 statt der Arbeitnehmerhaelfte}
\item[\texttt{quellensteuer}] \phantomsection\label{src:quellensteuer} Bundeszentralamt fuer Steuern (BZSt), 'Anrechenbarkeit der Quellensteuer auf Dividenden und Zinsen von Staaten, mit denen Deutschland ein Doppelbesteuerungsabkommen abgeschlossen hat', Stand 1. Januar 2025, Redaktionsschluss Juni 2025, laenderspezifische Zeilen, S. 1-25 (Belgien S.5, Daenemark S.7, Finnland S.9, Frankreich S.9, Irland S.12, Italien S.13, Kanada S.14, Niederlande S.21, Norwegen S.22, Spanien S.24, Schweiz S.23, Vereinigte Staaten S.24, Vereinigtes Koenigreich S.24). \url{https://www.bzst.de/SharedDocs/Downloads/DE/EU_OECD/anrechenbare_ausl_quellensteuer_2025.pdf}, consulté le 29.09.2026 (source primaire). \textit{Remarque de collecte~: PDF gespeichert unter refs/bzst-anrechenbare-quellensteuer-2025.pdf. Stand 1.1.2025, da 2026 zum Abrufzeitpunkt noch nicht veroeffentlicht war; DBA-Saetze aendern sich nur bei Vertragsaenderungen, siehe LUECKEN.md. DE-Zeile ist keine BZSt-Zeile, sondern aus abgeltungsteuer\_satz + soli\_satz abgeleitet.}
\item[\texttt{regelaltersgrenze}] \phantomsection\label{src:regelaltersgrenze} §35 SGB VI, Regelaltersgrenze fuer Jahrgaenge ab 1964. \url{https://www.deutsche-rentenversicherung.de/DRV/DE/Experten/Arbeitgeber-und-Steuerberater/summa-summarum/Lexikon/R/regelaltersgrenze}, consulté le 29.09.2026 (source primaire).
\item[\texttt{rentenwert\_2026}] \phantomsection\label{src:rentenwert_2026} Deutsche Rentenversicherung, Pressemitteilung 'Bundeskabinett beschliesst Rentenanpassung 2026' (29.04.2026), Anstieg von 40,79 auf 42,52 EUR. \url{https://www.deutsche-rentenversicherung.de/DRV/DE/Ueber-uns-und-Presse/Presse/Meldungen/2026/260429-rentenanpassung-2026-bundeskabinett.html}, consulté le 29.09.2026 (source primaire).
\item[\texttt{shiller\_url}] \phantomsection\label{src:shiller_url} Robert J. Shiller, Yale University, Online-Datentabelle S\&P 500 (Kurse, Dividenden, Ertraege, CPI) seit 1871. \url{http://www.econ.yale.edu/~shiller/data/ie_data.xls}, consulté le 29.09.2026 (source primaire). \textit{Remarque de collecte~: Erreichbarkeit per curl geprueft: HTTP 200, gueltige Excel-Datei (zuletzt gespeichert 2023, Autor RShiller)}
\item[\texttt{soli\_freigrenze\_2026}] \phantomsection\label{src:soli_freigrenze_2026} §3 Abs. 3 SolzG 1995 n.F. (Steuerfortentwicklungsgesetz 2024), erstmals VZ 2026. \url{https://www.buzer.de/gesetz/282/al234205-0.htm}, consulté le 29.09.2026 (source primaire).
\item[\texttt{soli\_satz}] \phantomsection\label{src:soli_satz} §4 SolzG 1995. \url{https://www.gesetze-im-internet.de/solzg_1995/__4.html}, consulté le 29.09.2026 (source primaire).
\item[\texttt{sparerpauschbetrag}] \phantomsection\label{src:sparerpauschbetrag} §20 Abs. 9 EStG, seit 2023. \url{https://www.gesetze-im-internet.de/estg/__20.html}, consulté le 29.09.2026 (source primaire). \textit{Remarque de collecte~: 2000 EUR bei Zusammenveranlagung}
\item[\texttt{sv\_arbeitnehmer}] \phantomsection\label{src:sv_arbeitnehmer} GKV-Spitzenverband-Faktenblatt 'Rechengroessen und Grenzwerte im Versicherungs- und Beitragsrecht fuer das Jahr 2026', S. 1. \url{https://www.gkv-spitzenverband.de/media/dokumente/presse/zahlen_und_grafiken/20260101_Faktenblatt_Rechengroessen_Beitragsrecht.pdf}, consulté le 29.09.2026 (source primaire). \textit{Remarque de collecte~: rv und av sind bereits die Arbeitnehmerhaelfte (halber Beitragssatz 18,6 \% bzw. 2,6 \%). kv\_allgemein und pv sind die vollen, hälftig geteilten Beitragssaetze (14,6 \% bzw. 3,6 \%); der Kinderlosenzuschlag von 0,6 \% traegt gesetzlich allein der Arbeitnehmer. kv\_zusatz (2,9 \%, voller Satz) ist eine Kopie von kv\_zusatzbeitrag\_2026, ebenfalls hälftig zu teilen (sv\_anteil in salaire.py teilt kv\_allgemein+kv\_zusatz durch 2).}
\item[\texttt{teilfreistellung\_aktienfonds}] \phantomsection\label{src:teilfreistellung_aktienfonds} §20 InvStG (Aktienfonds, natürliche Person im Privatvermögen). \url{https://www.gesetze-im-internet.de/invstg_2018/__20.html}, consulté le 29.09.2026 (source primaire).
\end{description}
